% TODO make hw stylesheet
\documentclass[11pt]{scrartcl}
\usepackage[a4paper, total={6in, 8in}]{geometry}
\usepackage[usenames,dvipsnames,svgnames]{xcolor}
\usepackage[shortlabels]{enumitem}
\usepackage[framemethod=TikZ]{mdframed}
\usepackage{amsmath,amssymb,amsthm}
\usepackage[colorlinks]{hyperref}
\usepackage{multicol}
\usepackage{tabularx}

\hypersetup{
	colorlinks=true,
	linkcolor=blue,
	filecolor=magenta,      
	urlcolor=magenta,
	pdftitle={MAT 319 HW}
}

\usepackage{microtype}
\usepackage{mathtools}
\usepackage[headsepline]{scrlayer-scrpage}
\usepackage{thmtools}
\usepackage{todonotes}

\mdfdefinestyle{mdgreenbox}{linecolor=ForestGreen,backgroundcolor=ForestGreen!5,
	linewidth=2pt,rightline=false,leftline=true,topline=false,bottomline=false,}
\declaretheoremstyle[headfont=\bfseries\sffamily\color{ForestGreen!70!black},
mdframed={style=mdgreenbox},headpunct={ --- },]{thmgreenbox}

\declaretheorem[style=thmgreenbox,name=Claim,numbered=no]{claim*}
\declaretheorem[style=thmgreenbox,name=Lemma,numbered=no]{lemma*}

\addtolength{\textheight}{3.14cm}
\providecommand{\ol}{\overline}
\providecommand{\ul}{\underline}
\providecommand{\eps}{\varepsilon}
\providecommand{\half}{\frac{1}{2}}
\providecommand{\dang}{\measuredangle} %% Directed angle
\providecommand{\CC}{\mathbb C}
\providecommand{\FF}{\mathbb F}
\providecommand{\NN}{\mathbb N}
\providecommand{\QQ}{\mathbb Q}
\providecommand{\RR}{\mathbb R}
\providecommand{\ZZ}{\mathbb Z}
\providecommand{\dg}{^\circ}
\providecommand{\ii}{\item}
\providecommand{\alert}[1]{{\sffamily\textbf{\textcolor{blue}{#1}}}}

\reversemarginpar

\theoremstyle{problem}
\newtheorem{innercustomthm}{Problem}
\newenvironment{problem}[1]
{\renewcommand\theinnercustomthm{#1}\innercustomthm}
{\endinnercustomthm}

\theoremstyle{remark}
\newtheorem{remarknumbered}{Remark}
\newenvironment{remark}
{\renewcommand\theremarknumbered{\hspace{-\the\fontdimen2\font\space}}\remarknumbered}
{\endremarknumbered}

\newenvironment{soln}{\begin{proof}[Solution]}{\end{proof}}
\newenvironment{walkthrough}{\noindent \textbf{Walkthrough.}}{\medskip}
\newlist{walk}{enumerate}{3}
\setlist[walk]{label=\bfseries (\alph*)}

\providecommand{\pow}{\operatorname{Pow}}
\providecommand{\vocab}[1]{{\textbf{\textcolor{ForestGreen}{#1}}}}
\providecommand{\lcm}{\operatorname{lcm}}
\providecommand{\abs}[1]{\left|#1\right|}

\usepackage{asymptote}
\begin{asydef}
	defaultpen(fontsize(10pt));
	unitsize(1cm);
	size(8cm); // set a reasonable default
	usepackage("amsmath");
	usepackage("amssymb");
	settings.tex="pdflatex";
	settings.outformat="pdf";
	// Replacement for olympiad+cse5 which is not standard
	import geometry;
	// recalibrate fill and filldraw for conics
	void filldraw(picture pic = currentpicture, conic g, pen fillpen=defaultpen, pen drawpen=defaultpen)
	{ filldraw(pic, (path) g, fillpen, drawpen); }
	void fill(picture pic = currentpicture, conic g, pen p=defaultpen)
	{ filldraw(pic, (path) g, p); }
	// some geometry
	pair foot(pair P, pair A, pair B) { return foot(triangle(A,B,P).VC); }
	pair orthocenter(pair A, pair B, pair C) { return orthocentercenter(A,B,C); }
	pair centroid(pair A, pair B, pair C) { return (A+B+C)/3; }
	// cse5 abbreviations
	path CP(pair P, pair A) { return circle(P, abs(A-P)); }
	path CR(pair P, real r) { return circle(P, r); }
	pair IP(path p, path q) { return intersectionpoints(p,q)[0]; }
	pair OP(path p, path q) { return intersectionpoints(p,q)[1]; }
	path Line(pair A, pair B, real a=0.6, real b=a) { return (a*(A-B)+A)--(b*(B-A)+B); }
	// cse5 more useful functions
	picture CC() {
		picture p=rotate(0)*currentpicture;
		currentpicture.erase();
		return p;
	}
	pair MP(Label s, pair A, pair B = plain.S, pen p = defaultpen) {
		Label L = s;
		L.s = "$"+s.s+"$";
		label(L, A, B, p);
		return A;
	}
	pair Drawing(Label s = "", pair A, pair B = plain.S, pen p = defaultpen) {
		dot(MP(s, A, B, p), p);
		return A;
	}
	path Drawing(path g, pen p = defaultpen, arrowbar ar = None) {
		draw(g, p, ar);
		return g;
	}
\end{asydef}

\usepackage{mathrsfs}

\ihead{\footnotesize MAT 319 HW04}
\ohead{\footnotesize September 25, 2024}

\title{\vspace{-2em}MAT 319 HW04}
\author{\large Aditya Pahuja}
\date{\large Due 9/25}
\begin{document}
\maketitle

\begin{problem}{12.1}
	Let $(s_n)$ and $(t_n)$ be sequences and suppose there exists $N_0$
	such that $s_n \le t_n$ for all $n > N_0$.
	Show $\liminf s_n \le \liminf t_n$ and $\limsup s_n \le \limsup t_n$.
\end{problem}

\begin{soln}
	Define the sets
	\[S_N = \{s_n : n > N\},\qquad
	T_N = \{t_n : n > N\}.\]

	\begin{claim*}
		For all $N > N_0$, $\inf S_N \le \inf T_N$ and $\sup S_N \le \sup T_N$.
	\end{claim*}
	\begin{proof}
		We know that $\inf S_N \le s_n$ for all $n > N$ by definition of $\inf$,
		so we have
		\[\inf S_N \le s_n \le t_n\]
		for all $n > N$ (noting $s_n \le t_n$ whenever $N > N_0$),
		which means that $\inf S_N$ is a lower bound for $T_N$.
		Since $\inf T_N$ is the greatest lower bound for $T_N$,
		$\inf S_N$ is at most as large as $\inf T_N$, which is what we wanted.
		
		The proof for $\sup S_N \le \sup T_N$ is analogous,
		but here it is anyway for completeness
		(yes, I copy-pasted and just changed the appropriate words).
		We know that $\sup T_N \ge t_n$ for all $n > N$ by definition of $\sup$,
		so we have
		\[\sup T_N \ge t_n \ge s_n\]
		for all $n > N$ (noting $t_n \ge s_n$ whenever $N > N_0$),
		which means that $\sup T_N$ is an upper bound for $S_N$.
		Since $\sup S_N$ is the least upper bound for $S_N$,
		$\sup T_N$ is at least as large as $\sup S_N$, which is what we wanted.
	\end{proof}

	Now, to finish, $\inf S_N \le \inf T_N$ being eventually true implies that
	\[\liminf s_n = \lim_{N\to\infty} \inf S_N \le \lim_{N\to\infty} \inf T_N
	= \liminf t_n,\]
	which gives $\liminf s_n \le \liminf t_n$.
	Similarly, because $\sup S_N \le \sup T_N$ eventually,
	\[\limsup s_n = \lim_{N\to\infty} \sup S_N \le \lim_{N\to\infty} \sup T_N
	= \limsup t_n,\]
	which gives $\limsup s_n \le \liminf t_n$.
\end{soln}

\pagebreak
\begin{problem}{12.3}
	Let $(s_n)$ and $(t_n)$ be the following sequences with period $4$:
	\begin{align*}
		(s_n) &= (0, 1, 2, 1, 0, 1, 2, 1, 0, 1, 2, 1, 0, 1, 2, 1, 0, \dots)\\
		(t_n) &= (2, 1, 1, 0, 2, 1, 1, 0, 2, 1, 1, 0, 2, 1, 1, 0, 2, \dots)
	\end{align*}
	Find
	
	\begin{tabularx}{\linewidth}{XX}
		\emph{\textbf{(a)}} $\liminf s_n + \liminf t_n$, &
		\emph{\textbf{(b)}} $\liminf(s_n + t_n)$,\\
		\emph{\textbf{(c)}} $\liminf s_n + \limsup t_n$, &
		\emph{\textbf{(d)}} $\limsup(s_n + t_n)$,\\
		\emph{\textbf{(e)}} $\limsup s_n + \limsup t_n$, &
		\emph{\textbf{(f)}} $\liminf(s_nt_n)$,\\
		\emph{\textbf{(g)}} $\limsup(s_nt_n)$ & 
	\end{tabularx}
\end{problem}

\begin{soln}
	First, we compute
	\begin{align*}
		(s_n+t_n) &= (2, 2, 3, 1, 2, 2, 3, 1, 2, 2, 3, 1, 2, \dots)\\
		(s_nt_n)  &= (0, 1, 2, 0, 0, 1, 2, 0, 0, 1, 2, 0, 0, \dots),
	\end{align*}
	noting that $(s_n+t_n)$ and $(s_nt_n)$ are both sequences with period $4$
	(namely $s_{n+4} = s_n$ and $t_{n+4} = t_n$ implies that
	$s_{n+4} + t_{n+4} = s_n + t_n$ and $s_{n+4}t_{n+4} = s_nt_n$
	for all $n$).
	
	The following claim about $\liminf$ and $\limsup$ of a periodic sequence
	will now completely solve the entire problem.
	\begin{claim*}
		Let $(x_n)$ be a periodic sequence with period $p$, meaning
		$x_n = x_{n+p}$ for all natural $n$. Then,
		\[\limsup x_n = \max\{x_1, x_2, \dots, x_p\},\quad
		\liminf x_n = \min\{x_1, x_2, \dots, x_p\}.\]
	\end{claim*}
	\begin{proof}
		Define $X_N = \{x_n : n > N\}$ and $P = \{x_1, x_2, \dots, x_p\}$.
		Since each element of $\{x_n\}$ is equal to one of $x_1$, \dots, $x_p$,
		$X_N\subseteq P$. On the other hand,
		picking an integer $q$ such that $pq > N$
		(which exists by the archimedean property of $\RR$)
		shows us that
		\[P = \{x_{pq + 1}, x_{pq + 2}, \dots, x_{pq + p}\} \subseteq X_N,\]
		which means we can conclude the equality of sets
		$X_N = P$.
		In particular,
		\[
		\sup X_N = \sup P = \max P,\qquad
		\inf X_N = \min P = \min P.
		\]
		because $\{x_1, \dots, x_p\}$ is finite.
		Taking the limit then finishes:
		\begin{align*}
			\limsup x_n &= \lim_{N\to\infty}\sup X_N
			= \lim_{N\to\infty} \max P = \max P,\\
			\liminf x_n &= \lim_{N\to\infty}\inf X_N
			= \lim_{N\to\infty} \min P = \min P.
			\qedhere
		\end{align*}
	\end{proof}
	
	Now, $(s_n)$, $(t_n)$, $(s_n + t_n)$, and $(s_nt_n)$ are all periodic, so:
	\begin{align*}
		\liminf s_n = \min\{0, 1, 2, 1\} = 0,\quad&
		\limsup s_n = \max\{0, 1, 2, 1\} = 2,\\
		\liminf t_n = \min\{2, 1, 1, 0\} = 0,\quad&
		\limsup t_n = \max\{2, 1, 1, 0\} = 2,\\
		\liminf s_n+t_n = \min\{2, 2, 3, 1\} = 1,\quad&
		\limsup s_n+t_n = \max\{2, 2, 3, 1\} = 3,\\
		\liminf s_nt_n = \min\{0, 1, 2, 0\} = 0,\quad&
		\limsup s_nt_n = \max\{0, 1, 2, 0\} = 2.
	\end{align*}
	
	Then, the answers are:
	
	\begin{tabularx}{\linewidth}{XX}
		{\textbf{(a)}} $\liminf s_n + \liminf t_n = 0 + 0 = 0$, &
		{\textbf{(b)}} $\liminf(s_n + t_n) = 1$,\\
		{\textbf{(c)}} $\liminf s_n + \limsup t_n = 0 + 2 = 2$, &
		{\textbf{(d)}} $\limsup(s_n + t_n) = 3$,\\
		{\textbf{(e)}} $\limsup s_n + \limsup t_n = 2 + 2 = 4$, &
		{\textbf{(f)}} $\liminf(s_nt_n) = 0$,\\
		{\textbf{(g)}} $\limsup(s_nt_n) = 2$ &
	\end{tabularx}
\end{soln}


\begin{problem}{12.6}
	Let $(s_n)$ be a bounded sequence, and let $k$ be a nonnegative real number.
	\begin{enumerate}[\textbf{(\alph*)}]
		\ii Prove $\limsup(ks_n) = k\cdot \limsup s_n$.
		\ii Do the same for $\liminf$.
		\ii What happens in (a) and (b) if $k < 0$?
	\end{enumerate}
\end{problem}

\begin{soln}
	Let $S$ be the set of subsequential limits of $(s_n)$
	and $kS$ be the set of subsequential limits of $(ks_n)$,
	the sequence obtained by multiplying each element of $s_n$ by $k$.
	(We use $kS$ to mean the set obtained by
	multiplying every element of $S$ with $k$.)
	All parts of the question will rely on the fact that
	\[kS = \{k\cdot\ell : \ell\in S\}.\]
	This is obvious when $k = 0$ because $(ks_{n_i})$ is the zero sequence,
	and so $S_0 = \{0\}$.
	When $k \ne 0$, we rephrase the set equality above as follows:
	\begin{claim*}
		The real number $\ell$ is in $S$ if and only if
		$k\cdot\ell$ is in $kS$.
	\end{claim*}
	\begin{proof}
		First, we show that, if $\ell$ is in $S$, then
		$k\cdot\ell$ is in $kS$.
		For any $\ell\in S$, there exists
		a subsequence $(s_{n_i})$ of $(s_n)$ converging to $\ell$.
		We know that
		\[\lim_{i\to\infty} ks_{n_i} = k\lim_{i\to\infty} s_{n_i}
		= k\cdot\ell.\]
		Obviously $(ks_{n_i})$ is a subsequence of the sequence $(ks_n)$;
		to be more thorough, we can let $t_n = ks_n$ and then say that
		the $t_{n_i}$ form a subsequence of $(t_n)$ because
		$(s_{n_i})$ being a subsequence of $(s_n)$
		implies that $n_i$ is an increasing function of $i$.
		Thus, $k\cdot\ell$ is a subsequential limit
		and therefore an element of $kS$.
		
		Now, we prove the converse, namely that
		$k\cdot\ell \in kS$ implies $\ell\in S$.
		In fact, we can just reuse the previous part:
		Let $k' = \frac1k$ (noting that $k\ne 0$)
		and define $t_n = ks_n$, so that $s_n = \frac{t_n}k = k't_n$.
		If some $k\cdot\ell$ is a subsequential limit of $(t_n)$
		(and thus in $kS$),
		then we showed above that $k'\cdot (k\cdot\ell) = \ell$
		must be a subsequential limit of the sequence $(k't_n) = (s_n)$,
		so $\ell\in S$.
	\end{proof}

	Now, we can solve (a) and (b) by making use of the fact that
	$\limsup(ks_n) = \sup kS$ and $\liminf(ks_n) = \inf kS$, since,
	when $k\ge 0$,
	\[\sup kS = k\sup S = k\cdot\limsup s_n,\qquad
	\inf kS = k\inf S = k\cdot\liminf s_n.\]
	For (c), we just have to recall that
	\[\liminf(s_n) = -\limsup(-s_n),\qquad \limsup s_n = -\liminf(-s_n),\]
	as this lets us get
	\[\liminf(ks_n) = -\limsup((-k)s_n) = -(-k)\cdot \limsup(s_n)
	= k\limsup s_n\]
	for $\liminf$, and
	\[\limsup(ks_n) = -\liminf((-k)s_n) = -(-k)\cdot \liminf(s_n)
	= k\liminf s_n\]
	for $\limsup$.
\end{soln}

\begin{problem}{12.8}
	Let $(s_n)$ and $(t_n)$ be bounded sequences of nonnegative numbers.
	Prove
	\[\limsup s_nt_n\le (\limsup s_n)(\limsup t_n).\]
\end{problem}

\begin{soln}
	Let $\ol s_N = \sup\{s_n : n > N\}$.
	Since the $s_n$ and $t_n$ are all nonnegative,
	$s_n \le \ol s_N$ for $n > N$ implies that
	\[s_nt_n \le \ol s_N t_n\]
	whenever $n > N$.
	This inequality is preserved when taking the $\limsup$
	(see Problem 12.1; here $N_0 = N$) so
	\[\limsup s_nt_n \le \ol s_N\cdot \limsup t_n\]
	for all $N$; note that we can factor out $\ol s_N$
	due to the previous problem (12.6).
	Then, we also take the limit as $N$ goes to infinity to say that
	\[\limsup s_nt_n = \lim_{N\to\infty} \limsup s_nt_n
	\le (\lim_{N\to\infty} \ol s_N)(\limsup t_n) = (\limsup s_n)(\limsup t_n),\]
	which is what we wanted to show.
\end{soln}

\begin{problem}{14.6}
	\begin{enumerate}[\textbf{(\alph*)}]
		\ii Prove that if $\sum |a_n|$ converges and $(b_n)$ is bounded
		then $\sum a_nb_n$ converges.
		\ii Observe that Corollary 14.7 is a special case of part (a).
	\end{enumerate}
\end{problem}

\begin{soln}
	For (a), it is equivalent to show that the sequence of partial sums
	of $a_nb_n$ is a Cauchy sequence. That is, for each $\eps > 0$,
	we want to show the existence of an $N$ such that $m,n > N$ implies that
	\[\abs{\sum_{k=m}^n a_kb_k} < \eps.\]
	Since $(b_n)$ is bounded, there exists an $M$ for which $|b_k| < M$
	for all $k$. Thus, we can get the following bound via triangle inequality:
	\[\abs{\sum_{k=m}^n a_kb_k} \le \sum_{k=m}^n\abs{a_kb_k}
	< \sum_{k=m}^nM\abs{a_k} = M\sum_{k=m}^n \abs{a_k}.\]
	Now, since $\sum|a_n|$ converges and $\frac\eps M$ is a positive real,
	there exists an $\hat N$ such that, whenever $m, n > \hat N$,
	\[\sum_{k = m}^n |a_n| \le \frac{\eps}{M},\]
	which means that $m, n > \hat N$ implies
	\[\abs{\sum_{k=m}^n a_kb_k}
	< M\sum_{k=m}^n\abs{a_k} < \frac\eps{M}\cdot M = \eps.\]
	This is enough to show that the partial sums of $a_kb_k$
	form a Cauchy sequence (taking $N = \hat N$), thus $\sum a_nb_n$ converges.
	
	As for part (b), Corollary 14.7 (absolute convergence implies convergence)
	arises from taking $b_n = 1$ for all $n$.
\end{soln}

\pagebreak
\begin{problem}{14.7}
	Prove that if $\sum a_n$ is a convergent series of nonnegative numbers
	and $p > 1$, then $\sum a_n^p$ converges.
\end{problem}

\begin{soln}
	Since $\sum a_n$ converges, $\lim_{n\to\infty} a_n = 0$.
	This means that eventually $a_n < 1$, say for $n > N$
	where $N$ is an integer
	(to be precise, taking $\eps = 1$, there exists an $N$
	such that $n > N$ implies $|a_n - 0| = a_n < \eps$).
	For $n > N$, we can then say that $|a_n^p| = a_n^p < a_n$
	so, by the comparison test,
	\[\sum_{n = N+1}^\infty a_n^p \le \sum_{n=N+1}^\infty a_n,\]
	in particular the left-hand sum converges.
	This is equivalent to the statement
	\begin{equation}
		\text{for all $\eps > 0$, there exists a $B$ such that }
		n > m > B \implies \abs{\sum_{k=m}^n a_k^p} < \eps
	\end{equation}
	by the Cauchy criterion for the series $\sum_{k > N} a_k^p$,
	where the sum is taken over all indices greater than $N$.
	On the other hand, this is also exactly the Cauchy criterion
	for $\sum a_k^p$ just taken over all natural numbers:
	that is, the chain of logic looks like
	\[\sum_{k=N+1}^\infty a_k^p\text{ converges} \iff (1) \iff
	\sum_{k=1}^\infty a_k^p \text{ converges},\]
	so $\sum a_k^p$ converges!
\end{soln}

\begin{problem}{14.9}
	For sequences $(a_n)$ and $(b_n)$, suppose that the set
	$\{n\in\NN : a_n\ne b_n\}$ is finite. Show that $\sum a_n$
	converges if and only if $\sum b_n$ converges.
\end{problem}

\begin{soln}
	We show that $\sum a_n$ satisfies the Cauchy criterion if and only if
	$\sum b_n$ also does.
	
	Let $n_0$ be an integer larger than every element of
	$\{n\in\NN: a_n\ne b_n\}$, so that $k > n_0$ implies $a_k = b_k$.
	Then, $q > p > n_0$ implies that $a_k = b_k$ if $p\le k\le q$,
	hence $\sum_{k=p}^q a_k = \sum_{k=p}^q b_k$.
	From there we can see that
	\begin{equation}
		\text{for all $\eps > 0$, there exists an $N$ such that }
		q > p > \max\{N, n_0\} \implies \abs{\sum_{k=p}^q a_k} < \eps
	\end{equation}
	is the exact same statement as
	\begin{equation}
		\text{for all $\eps > 0$, there exists an $N$ such that }
		q > p > \max\{N, n_0\} \implies \abs{\sum_{k=p}^q b_k} < \eps
	\end{equation}
	The former statement is the Cauchy criterion for $\sum a_k$
	and the latter is the Cauchy criterion for $\sum b_k$,
	so we have
	\[\sum_{k=1}^\infty a_k\text{ converges} \iff (1) \iff (2) \iff
	\sum_{k=1}^\infty b_k\text{ converges},\]
	which is what we wanted.
\end{soln}

\pagebreak
\begin{problem}{14.14}
	Prove that $\sum \frac1n$ diverges by comparison with
	$\sum_{n\ge 2} a_n$ where $(a_n)$ is the sequence
	\[\left(\frac12, \frac14, \frac14, \frac18, \frac18, \frac18, \frac18,
	\frac1{16}, \frac1{16}, \frac1{16}, \frac1{16}, \frac1{16}, \frac1{16},
	\frac1{16}, \frac1{16}, \frac1{32}, \frac1{32},\dots\right).\]
	In general, $a_n = \frac{1}{2^r}$, where $r$ is the integer
	satisfying $2^{r-1} < n \le 2^r$
	(and the sequence is indexed starting at $2$ for some ridiculous reason).
\end{problem}

\begin{soln}
	We start by immediately correcting the horrendous indexing given to us:
	define the sequence $(b_n)$ over all naturals such that
	$b_n = a_{n+1}$ for all $n$. Our goal is to prove the divergence
	of the harmonic series via comparison with $\sum b_n$.
	
	\begin{claim*}
		Let $(s_n)$ be the sequence of partial sums of $(b_n)$.
		Then, for every natural number $k$,
		\[s_{2^k - 1} = \frac k2.\]
	\end{claim*}
	\begin{proof}
		We induct on $k$. The base case $k = 1$ is trivial, as
		\[s_1 = a_1 = \frac 12.\]
		
		Supposing the claim holds for some $k$ (so $s_{2^k - 1} = \frac k2$),
		we can write
		\[s_{2^{k+1} - 1}
		= s_{2^k - 1} + \sum_{j=2^k}^{2^{k+1} - 1} a_j
		= \frac k2 + \sum_{j=2^k}^{2^{k+1} - 1} \frac{1}{2^{k+1}}
		= \frac{k}{2} + 2^k\cdot \frac{1}{2^{k+1}}
		= \frac k2 + \frac 12 = \frac{k+1}{2},\]
		noting that $a_j = \frac{1}{2^{k+1}}$ when $2^k \le j < 2^{k+1}$
		because the smallest power of two larger than $j$ is $2^{k+1}$.
		Thus, the claim is proven by the principle of mathematical induction.
	\end{proof}
	
	Now, since each $b_n$ is positive, $s_{n+1} - s_n = a_{n+1} > 0$
	means that $(s_n)$ is monotonically increasing,
	so $\lim_{n\to\infty} s_n$ exists. Moreover, any subsequence of $(s_n)$
	converges to this same limit. Taking $n_k = 2^k - 1$,
	we look at the limit of the subsequence of $(s_n)$ given by $(s_{n_k})$:
	\[\sum_{n=1}^\infty b_n = \lim_{n\to\infty} s_n = \lim_{k\to\infty} s_{n_k}
	= \lim_{k\to\infty} \frac k2 = +\infty.\]
	Thus, $\sum b_n$ diverges to $+\infty$.
	
	Now, by definition, for any natural number $n$, we know that
	$a_{n+1} = \frac{1}{2^r}$ where $r$ is
	an integer such that $2^{r-1} < n + 1 \le 2^r$.
	Thus, $n < 2^r$ implies
	\[\frac{1}{n} > \frac{1}{2^r} = a_{n+1} = b_n.\]
	Since $(b_n)$ is a sequence of nonnegative numbers
	and $\frac 1n > b_n$ for all $n$,
	the comparison theorem tells us that $\sum b_n$ diverging to
	$+\infty$ means that $\sum \frac 1n$ also diverges to $+\infty$.
\end{soln}


\end{document}